\section*{الفصل \ref{ch:g5:air-pressure-weather} --- الضغط الجوي والطقس}

\begin{solution}{exo:g5:air-pressure-weather:1}
تُوزن الكرة نفسها منفوخة ثم مفرَّغة: ولم يغادرها إلا \omterm{def:g2:air-around-us:air}{الهواء}،
فيقرأ الميزان أقلّ --- والفرق \omterm{def:g4:weight-and-mass:weight}{وزن} \omterm{def:g2:air-around-us:air}{الهواء} الهارب. \omterm{def:g2:air-around-us:air}{فالهواء} مادة
لها \omterm{def:g3:measuring-mass:mass}{كتلة}.
\end{solution}

\begin{solution}{exo:g5:air-pressure-weather:2}
هو دفع \omterm{def:g2:air-around-us:air}{الهواء} المعصور في قاع \omterm{def:g5:air-pressure-weather:pressure}{الغلاف الجوي} على كل سطح يلمسه.
وهو يأتي من \omterm{def:g4:weight-and-mass:weight}{وزن} محيط \omterm{def:g2:air-around-us:air}{الهواء} العميق الذي فوقنا، ويدفع من كل جهة
--- إلى الأسفل وجانبًا وإلى الأعلى.
\end{solution}

\begin{solution}{exo:g5:air-pressure-weather:3}
لأن الدفعات تأتي متعادلة: \omterm{def:g2:air-around-us:air}{فالهواء} في الجهة الأخرى من راحتك ---
وداخل جسمك --- يدفع بالشدّة نفسها بالضبط. ولا تُظهر جبروتَها
إلا دفعة \emph{غير متعادلة}، مثل التي على البطاقة أو على
القنّينة المتكرمشة.
\end{solution}

\begin{solution}{exo:g5:air-pressure-weather:4}
\omterm{def:g4:weight-and-mass:weight}{وزن} الماء يدفع البطاقة إلى الأسفل؛ \omterm{def:g5:air-pressure-weather:pressure}{والغلاف الجوي} يدفعها إلى
الأعلى. ويفوز \omterm{def:g5:air-pressure-weather:pressure}{الغلاف الجوي} --- فتبقى البطاقة ويُحبس الماء.
\end{solution}

\begin{solution}{exo:g5:air-pressure-weather:5}
أنت تنزع \emph{\omterm{def:g2:air-around-us:air}{الهواء}} من داخل الماصّة. فيدفع \omterm{def:g5:air-pressure-weather:pressure}{الغلاف الجوي}،
الضاغط على العصير في الكأس، العصيرَ صاعدًا في الماصّة المُخلاة:
فالسماء هي التي ترفع.
\end{solution}

\begin{solution}{exo:g5:air-pressure-weather:6}
ضغط الكأس مستوية عصر \omterm{def:g2:air-around-us:air}{الهواء} من خلفها. ومع انعدام ما يدفع من
تحتها يثبّت دفعُ \omterm{def:g5:air-pressure-weather:pressure}{الغلاف الجوي} الكأسَ على البلاط.
\end{solution}

\begin{solution}{exo:g5:air-pressure-weather:7}
يضعف --- إذ يبقى فوقك محيط \omterm{def:g2:air-around-us:air}{هواء} أقلّ ليضغط إلى الأسفل. ومن
علاماته: طقطقة الأذنين في صعود سريع، وضيق النفس في \omterm{def:g2:air-around-us:air}{الهواء}
الرقيق (وشاي يغلي أبرد ممّا ينبغي).
\end{solution}

\begin{solution}{exo:g5:air-pressure-weather:8}
دفع \omterm{def:g5:air-pressure-weather:pressure}{الغلاف الجوي} --- أي \omterm{def:g5:air-pressure-weather:pressure}{الضغط الجوي}. والهبوط السريع ينذر بعاصفة
في الطريق: فانقل النزهة إلى الداخل.
\end{solution}

\begin{solution}{exo:g5:air-pressure-weather:9}
أُغلقت القنّينة على القمّة فحبست \omterm{def:g2:air-around-us:air}{هواء} جبل رقيقًا في داخلها؛
وفي الوادي صار الدفع الخارجي أشدّ من الدفع المحبوس في الداخل،
فانطوت القنّينة. أما في الاتجاه الآخر فقنّينة أُغلقت في الوادي
تحفظ هواءً \emph{قويًّا} في داخلها بينما يضعف الدفع الخارجي مع
الارتفاع: فتنتفخ مشدودة --- وكيس الرقائق المغلق يفعل مثل ذلك في
كل رحلة جبلية.
\end{solution}

\begin{solution}{exo:g5:air-pressure-weather:10}
ينسكب \omterm{def:g2:air-around-us:air}{الهواء} دائمًا من مناطق الضغط المرتفع الثقيلة نحو مناطق
الضغط المنخفض الخفيفة، وذلك الانسكاب هو \omterm{def:g2:air-around-us:wind}{الريح}. فيكتمل التعريف
القديم بنفسه: \omterm{def:g2:air-around-us:wind}{الريح} \omterm{def:g2:air-around-us:air}{هواء} متحرّك --- حرّكته فروق \omterm{def:g5:air-pressure-weather:pressure}{الضغط الجوي}.
\end{solution}

\begin{solution}{exo:g5:air-pressure-weather:11}
مع الارتفاع يضعف دفع \omterm{def:g5:air-pressure-weather:pressure}{الغلاف الجوي} على الماء، فيأتي الغليان
أسهل: فتنزلق \omterm{def:g4:changes-of-state:boiling-point}{درجة الغليان} دون \qty{100}{\celsius} بكثير. وبقانون
المسطبة لا يستطيع الماء الغالي أن يصعد فوق درجة غليانه أبدًا ---
واللهب الزائد لا يصنع إلا بخارًا. فماء القمّة يغلي باردًا
\emph{ويبقى هناك}: فالشاي ضعيف، والمعكرونة قاسية، والفيزياء
راضية.
\end{solution}
